\section{Experimental designs that isolate context use}\label{sec:C11:experiments}
A proposed long-context evaluation should begin with matched information access. Full-history, retrieval, and summary systems should receive the same underlying evidence, tools, model version, and authorization. The comparison then varies only how evidence reaches the model. A second experiment can permit learned context control, but its training and search costs should be reported separately. Otherwise, architectural selection effects become entangled with an additional optimization budget.

Recent work makes these distinctions concrete. The Compaction Cliff classifies knowledge by type and applies separate operators for compaction, decomposition, and retrieval, retaining recognized constraints under type-specific rules~\cite{zerhoudi2026cliff}. Its guarantees depend on constraint classification, applicable scope, feasible budgets, and verification; they are not unconditional preservation of every real safety requirement. Its behavioral retail comparison is not token-matched, and some reference labels are model-generated. Those qualifications prevent an observed advantage from being assigned solely to typing. The preprint was available before the cutoff even though its listed conference dates fall later.

The methodological conflict is useful. Uniform summarization assumes that one semantic compression objective is adequate across records, whereas typed retention treats some distinctions as constraints that should survive with stronger fidelity. The latter can preserve a rare prohibition even when a relevance score prefers repeated descriptive content. It can also fail if the classifier mistakes a declarative fact for ordinary background or assigns an inappropriate scope. A test of constraint preservation should therefore vary grammatical form and authority independently of content frequency.

Two 2026 surveys supply organizing lenses rather than independent experimental confirmation. The rate--distortion view groups several forms of context-derived compression under budget and downstream-utility tradeoffs~\cite{colaco2026ratedistortion}. The October selective-memory survey distinguishes information participating in computation from stored information that can later be activated~\cite{gou2026selective}. These taxonomies help align accounting boundaries, but they do not prove that cache eviction, semantic forgetting, and text summarization have identical guarantees. Their state transitions and recovery mechanisms remain different.

An original deduction follows from combining these distinctions with the sufficiency criterion: fidelity should be allocated according to the loss of the distinction being removed, while activation cost should be charged when stored evidence becomes computation. This does not yield a universal scalar salience score. It yields a falsifiable comparison: preserve the same constraint set, vary non-constraint evidence and activation policy, and measure both justified action and full resource use. A method that wins only by silently relaxing the constraint set has changed the task.

Use at least three baselines: a full-history baseline within the admissible budget, a recency window, and a simple retrieval or structured-state baseline. The full-history baseline tests whether the proposed transformation removes harmful distraction or merely compensates for capacity. The recency baseline measures whether most tasks are effectively short. The simple structured baseline tests whether elaborate learned control is necessary. An oracle evidence condition is useful as a diagnostic ceiling, provided it is clearly labeled and never compared as an executable production method.

For each task, instrument the first decision at which a required historical distinction matters. Record whether the distinction exists in storage, appears in the model input, is interpreted correctly, and leads to the correct tool action. This decomposition prevents a retrieval miss from being mistaken for weak reasoning and prevents a correct answer with an invalid side effect from being counted as success. The same decomposition supports targeted ablations: preserve the record but perturb its timestamp, preserve its content but remove its source role, or preserve both while moving its position.

Uncertainty should reflect the unit of generalization. Repeated decoding seeds on the same task measure stochastic policy variation, not variation across task families. Use task-level resampling for task uncertainty, environment-level separation where tasks share state, and independent training seeds when comparing learned context policies. Task order matters whenever memory or a learned skill bank persists. A single favorable sequence can behave as an implicit curriculum and should not be treated as evidence of stable continual improvement.

Contamination testing should include withheld corrections and contradictory updates rather than only unseen questions. A system that remembers the original fact but fails to adopt an authoritative revision has failed a different requirement from one that never stored the fact. Include negative controls where an old record should be ignored and where no historical record is relevant. Without these controls, systems that overapply memory may appear strong on a benchmark constructed entirely from positive recall opportunities.

The proposed experiment is discriminating only if the result can falsify a mechanism. If a method claims to preserve decision-critical evidence, a loss of that evidence under matched storage and retrieval is directly informative. A small change in an aggregate score without an inspectable failure record may be insufficient to identify whether retention, attention, or action selection changed.

